\lesson{1}{Oct 11 2021 Mon (10:42:23)}{GCF and SP}{Unit 2}

\subsubsection{Difference of Squares Binomials}
\label{sub_sub_sec:difference_of_squares_binomials}

\begin{definition}[Difference of Squares Binomials]
  \label{def:difference_of_squares_binomials}

  A \gls{dsb} is when both terms in a binomial are perfect squares \textbf{AND}
  they are being subtracted. A \gls{dsb} can be factored using the following
  patter:

  \begin{equation}
    \label{eqt:dsb_pattern}
    \begin{split}
      a^{2} - b^{2} = (a + b)(a - b)
    \end{split}
  .\end{equation}
\end{definition}

\begin{note}
  \label{nte:dsb_pattern_note}

  The square root of each perfect square $(a^{2} \quad b^{2})$ is $a \quad b$.
  These terms are placed inside two sets of parentheses where one is addition
  and one is subtraction. Multiplying the two factors $(a + b)(a - b)$ results
  in the following:

  \begin{equation}
    \label{eqt:dsb_pattern_result}
    \begin{split}
      (a + b)&(a - b) \\
      a^{2} + ab &-ab - b^{2} \\
      a^{2} + 0&ab - b^{2} \\
      a^{2} &- b^{2} \\
    \end{split}
  .\end{equation}
\end{note}

Here's how you identify a \gls{dsb}.

\begin{enumerate}
  \label{enum:dsb_identification}

  \item What do you notice about the first terms in each of the products:
    $g^{2} - 16 \quad 25g^{2} - 64$?

    The terms $g^{2}$ and $25g^{2}$ are perfect squares. 

  \item What do you notice about the last terms in each of the products:
    $g^{2} - 16 \quad 25g^{2} - 64$?

    The terms $16$ and $64$ are also perfect squares.

  \item What sign separates the first and last terms in each of the products:
    $g^{2} - 16 \quad 25g^{2} - 64$?

    They are all subtraction signs.

  \item Now try working backwards. What two sets of binomials that multiply
    together to produce $g^{2} - 16$?

    The two binomials are: $g + 4 \quad g - 4$.

  \item Why are the signs different for each of the factors?

    The signs are different so that, when the two sets of parentheses are
    multiplied, the last term will be negative and the middle terms will cancel.
\end{enumerate}

\begin{example}
  \label{exm:dsb_identification}

  Check if $64d^{4} - 361d^{2}$ is a \gls{dsb}:

  \begin{equation}
    \label{eqt:dsb_identification}
    \begin{split}
      \textrm{Step 1:} \quad & \quad \textrm{Look at the first term} \\
        64d^{4} \quad & \quad \textrm{Is it a perfect square? Yes!} \quad \sqrt{64} = 8 \\
      \textrm{Step 2:} \quad & \quad \textrm{Look at the last term} \\
        361d^{4} \quad & \quad \textrm{Is it a perfect square? Yes!} \quad \sqrt{361} = 19 \\
      \textrm{Step 3:} \quad & \quad \textrm{What sign separates the first and last term} \\
        64d^{4} - 361d^{4} \quad & \quad \textrm{The $-$ sign, which is what's needed.} \\
      \textrm{Step 4:} \quad & \quad \textrm{Let's find the two sets of binomials!} \\
        & a = \sqrt{64d^{4}} = 8d^{2} \\
        & b = \sqrt{361d^{2}} = 19d \\
        & (8d^{2} + 19d)(8d^{2} - 19d) \\
        & \boxed{d^{2}(8d + 19d)(8d - 19d)} \quad \textrm{Factor the $d^{2}$} \\
    \end{split}
  .\end{equation}
\end{example}

% subsubsection difference_of_squares_binomials (end)

\subsubsection{Perfect Square Trinomials}
\label{perfect_square_trinomials}

\begin{definition}[Difference of Squares Trinomials]
  \label{def:difference_of_squares_trinomials}

  A \gls{pst} is when first and last term of the \gls{pst} are both perfect
  squares \textbf{AND} the middle term is made from two numbers that when you
  multiply you get the last term and when you add them, you get the middle term.

  Here's the pattern that a \gls{pst} follows:

  \begin{equation}
    \label{eqt:dst_pattern}
    \begin{split}
      a^{2} - 2ab + b^{2} = (a - b)^{2}
    \end{split}
  .\end{equation}
\end{definition}

Here's how you identify a \gls{pst}.

\begin{enumerate}
  \label{enum:dst_identification}

  \item What do you notice about the first terms in each of the products:
    $g^{2} + 12g + 36 \quad g^{2} - 6g + 9$ \\

    They are all perfect squares.

  \item What do you notice about the last terms in each of the products:
    $g^{2} + 12g + 36 \quad g^{2} - 6g + 9$ \\

    They too, are perfect squares.

  \item What signs do the first and last terms have?

    They are all positive.

  \item What about the sign on the middle term in each product?

    The sign of the middle term is identical to the sign in the factors
    (squared binomials).

  \item Now try working backwards. What two sets of binomials that multiply
    together to produce $g^{2} - 16g + 9$?

    The binomials are: $(g - 8)(g - 8)$ or $(g - 8)^{2}$.

  \item How is the middle term of $-16g$ used to find the factors?

    The middle term is used to determine the operation inside each set of
    parentheses and to verify this polynomial is a \gls{pst}. $-8$ and $g$ are
    multiplied to and doubled to arrive at $-16g$.
\end{enumerate}

\begin{example}
  \label{exm:dst_pattern}

  Let's find the factors of the following \gls{pst}: $x^{2} - 5x - 24$:

  \begin{equation}
    \label{eqt:dst_pattern_result}
    \begin{split}
      & x^{2} - 5x - 24 \\
      & 3 \times -8 = -24  \\
      & 3 + -8 = -5 \\
      & (x + a)(x + b) \\
      & \boxed{(x + 3)(x - 8)} \\
    \end{split}
  .\end{equation}

  I didn't just pluck $3$ and $-8$ out of nowhere. I tried a bunch of numbers by
  multiplying them together and seeing if they return $-24$. Then, I add them to
  see if they return $-5$. If so, those are the numbers that are needed.
\end{example}

% subsubsection difference_of_squares_trinomials (end)

\newpage
